\FormulaThmDeltaK[Die Identität ist ein Automorphismus]{%
 m\in\operatorname{Assoc}(S)\vdash\Id_S\in\operatorname{Aut}(S,m)
}{B37TAutomorphismIdentity}{\DeltaRow{Mengen}{S}\DeltaRow{Funktionen}{m}}
\begin{tabproofwide}
 \proofstepwidestar[1]{m\in\operatorname{Assoc}(S)}{\rA}
 \proofstepwidestar[]{\Id_S\colon S\bij S}{\FormulaRefAuto{\Id_A\colon A\bij A}[thm]}
 \proofstepwidestar[]{\Id_S\in\operatorname{Sym}(S)}{\FormulaRefAuto{FiniteSemigroupSymmetricSetDef}[def]{2}}
 \proofstepwidestar[1]{m^{\Id_S}=m}{\rAEa{\rAEb{\rAEb{\rAEb{\FormulaRefAuto{FiniteSemigroupRelabellingAction}[thm]{1,\rAI{3,3}}}}}}}
 \proofstepwidestar[1]{\Id_S\in\operatorname{Aut}(S,m)}{\FormulaRefAuto{FiniteSemigroupAutomorphismGroupDef}[def]{1,3,4}}
\end{tabproofwide}

\FormulaThmDeltaK[Automorphismen sind unter Komposition abgeschlossen]{%
 \begin{aligned}[t]
 &m\in\operatorname{Assoc}(S)\dsep\pi,\sigma\in\operatorname{Aut}(S,m)\vdash{}\\[-2pt]
 &\sigma\circ\pi\in\operatorname{Aut}(S,m)
 \end{aligned}
}{B37TAutomorphismComposition}{\DeltaRow{Mengen}{S}\DeltaRow{Funktionen}{m\dsep\pi\dsep\sigma}}
\begin{tabproofwide}
 \proofstepwidestar[1]{m\in\operatorname{Assoc}(S)}{\rA}
 \proofstepwidestar[2]{\pi\in\operatorname{Aut}(S,m)}{\rA}
 \proofstepwidestar[3]{\sigma\in\operatorname{Aut}(S,m)}{\rA}
 \proofstepwidestar[1,2]{\pi\in\operatorname{Sym}(S)\land m^\pi=m}{\FormulaRefAuto{FiniteSemigroupAutomorphismGroupDef}[def]{1,2}}
 \proofstepwidestar[1,3]{\sigma\in\operatorname{Sym}(S)\land m^\sigma=m}{\FormulaRefAuto{FiniteSemigroupAutomorphismGroupDef}[def]{1,3}}
 \proofstepwidestar[1,2,3]{\begin{aligned}[t]&\Id_S\in\operatorname{Sym}(S)\land\sigma\circ\pi\in\operatorname{Sym}(S)\\[-2pt]&{}\land m^\pi\in\operatorname{Assoc}(S)\land m^{\Id_S}=m\land(m^\pi)^\sigma=m^{\sigma\circ\pi}\end{aligned}}{\FormulaRefAuto{FiniteSemigroupRelabellingAction}[thm]{1,\rAI{\rAEa{4},\rAEa{5}}}}
 \proofstepwidestar[1,2,3]{\sigma\circ\pi\in\operatorname{Sym}(S)}{\rAEa{\rAEb{6}}}
 \proofstepwidestar[1,2,3]{(m^\pi)^\sigma=m^{\sigma\circ\pi}}{\rAEb{\rAEb{\rAEb{\rAEb{6}}}}}
 \proofstepwidestar[1,2,3]{m^\sigma=m^{\sigma\circ\pi}}{\rIE{\rAEb{4},8}}
 \proofstepwidestar[1,2,3]{m^{\sigma\circ\pi}=m^\sigma}{\FormulaRefAuto{a=b\vdash b=a}[thm]{9}}
 \proofstepwidestar[1,2,3]{m^{\sigma\circ\pi}=m}{\rChain{10,\rAEb{5}}}
 \proofstepwidestar[1,2,3]{\sigma\circ\pi\in\operatorname{Aut}(S,m)}{\FormulaRefAuto{FiniteSemigroupAutomorphismGroupDef}[def]{1,7,11}}
\end{tabproofwide}

\FormulaThmDeltaK[Das Inverse eines Automorphismus ist ein Automorphismus]{%
 m\in\operatorname{Assoc}(S)\dsep\pi\in\operatorname{Aut}(S,m)
 \vdash\pi^{-1}\in\operatorname{Aut}(S,m)
}{B37TAutomorphismInverse}{\DeltaRow{Mengen}{S}\DeltaRow{Funktionen}{m\dsep\pi}}
\begin{tabproofwide}
 \proofstepwidestar[1]{m\in\operatorname{Assoc}(S)}{\rA}
 \proofstepwidestar[2]{\pi\in\operatorname{Aut}(S,m)}{\rA}
 \proofstepwidestar[1,2]{\pi\in\operatorname{Sym}(S)\land m^\pi=m}{\FormulaRefAuto{FiniteSemigroupAutomorphismGroupDef}[def]{1,2}}
 \proofstepwidestar[1,2]{\pi\colon S\bij S}{\FormulaRefAuto{FiniteSemigroupSymmetricSetBijection}[thm]{\rAEa{3}}}
 \proofstepwidestar[1,2]{\pi^{-1}\colon S\bij S}{\FormulaRefAuto{InverseFunctionBijection}[thm]{4}}
 \proofstepwidestar[1,2]{\pi^{-1}\in\operatorname{Sym}(S)}{\FormulaRefAuto{FiniteSemigroupSymmetricSetDef}[def]{5}}
 \proofstepwidestar[1,2]{\pi^{-1}\circ\pi=\Id_S}{\FormulaRefAuto{InverseFunctionLeftComposition}[thm]{4}}
 \proofstepwidestar[1,2]{(m^\pi)^{\pi^{-1}}=m^{\pi^{-1}\circ\pi}}{\rAEb{\rAEb{\rAEb{\rAEb{\FormulaRefAuto{FiniteSemigroupRelabellingAction}[thm]{1,\rAI{\rAEa{3},6}}}}}}}
 \proofstepwidestar[1,2]{m^{\pi^{-1}}=m^{\Id_S}}{\rIE{\rAEb{3},7,8}}
 \proofstepwidestar[1,2]{m^{\Id_S}=m}{\rAEa{\rAEb{\rAEb{\rAEb{\FormulaRefAuto{FiniteSemigroupRelabellingAction}[thm]{1,\rAI{\rAEa{3},6}}}}}}}
 \proofstepwidestar[1,2]{m^{\pi^{-1}}=m}{\rChain{9,10}}
 \proofstepwidestar[1,2]{\pi^{-1}\in\operatorname{Aut}(S,m)}{\FormulaRefAuto{FiniteSemigroupAutomorphismGroupDef}[def]{1,6,11}}
\end{tabproofwide}
Die drei Tabellen beweisen die behauptete Stabilität unter Identität,
Komposition und Inversen. Assoziativität der Komposition und die
Identitäts- und Inversengesetze werden von den Selbstbijektionen geerbt;
eine Gruppendefinition aus einem späteren Band wird hier nicht vorausgesetzt.